% ---------- 多語言 ----------
\usepackage{polyglossia}
\providecommand{\ScriptsDebugFontRoot}{./fonts/}
\setdefaultlanguage{english}
\setmainfont{CMU Serif}

% 天城體
\setotherlanguage{sanskrit}
\newfontfamily\sanskritfont[Script=Devanagari]{TiroDevanagariSanskrit-Regular.ttf}[Path=\ScriptsDebugFontRoot]
\newcommand{\SA}[1]{\textsanskrit{#1}}

% 阿拉伯文
\setotherlanguage{arabic}
\newfontfamily\arabicfont[Script=Arabic]{NotoNaskhArabic-Regular.ttf}[Path=\ScriptsDebugFontRoot]
\newcommand{\AR}[1]{\textarabic{#1}}

% 注音符號
\newCJKfontfamily\bopomofofont{I.Ming-8.10.ttf}[Path=\ScriptsDebugFontRoot]
\newcommand{\ZY}[1]{{\bopomofofont #1}}

% 強制調用 times.ttf（不限定腳本/語言）
\newfontfamily\timesfont{times.ttf}[Path=\ScriptsDebugFontRoot]
\newcommand{\TIMES}[1]{{\timesfont #1}}

% 希伯來文
\setotherlanguage{hebrew}
\newfontfamily\hebrewfont[Script=Hebrew]{NotoSerifHebrew-Regular.ttf}[Path=\ScriptsDebugFontRoot]
\newcommand{\HE}[1]{\texthebrew{#1}}

% 藏文
\setotherlanguage{tibetan}
\newfontfamily\tibetanfont[Script=Tibetan]{NotoSerifTibetan-Regular.ttf}[Path=\ScriptsDebugFontRoot]
\newcommand{\TI}[1]{\texttibetan{#1}}

% 國際音標
\newfontfamily\ipafont{CMU Serif} % 或 CMUSerif-Roman
\newcommand{\IPA}[1]{{\ipafont #1}}

% 諺文及朝鮮漢字
\newCJKfontfamily\koreanfamily{NotoSerifKR-Regular.ttf}[Path=\ScriptsDebugFontRoot]
\newcommand{\KO}[1]{%
  {\xeCJKsetup{CJKspace=true}\koreanfamily #1}%
}

% 大陸字形
\newCJKfontfamily\SCfont{NotoSerifSC-Regular.ttf}[Path=\ScriptsDebugFontRoot]
\newcommand{\SC}[1]{{\SCfont #1}}

% 日文
\newCJKfontfamily\JPfont{NotoSerifJP-Regular.ttf}[Path=\ScriptsDebugFontRoot]
\newcommand{\JP}[1]{{\JPfont #1}}

% 臺灣字形
\newCJKfontfamily\TCfont{NotoSerifTC-Regular.ttf}[Path=\ScriptsDebugFontRoot]
\newcommand{\TC}[1]{{\TCfont #1}}

% 越南國語字
\newfontfamily\VIfont{Minh Nguyen Regular.ttf}[Path=\ScriptsDebugFontRoot]
\newcommand{\VI}[1]{{\VIfont #1}}

% 漢喃文
\newCJKfontfamily\NOMfont{Minh Nguyen Regular.ttf}[Path=\ScriptsDebugFontRoot]
\newcommand{\NOM}[1]{{\NOMfont #1}}

% 西夏文
\newCJKfontfamily\TGfont{NotoSerifTangut-Regular.ttf}[Path=\ScriptsDebugFontRoot]
\newcommand{\TG}[1]{{\TGfont #1}}

% 楔形文字
\newfontfamily\CUfont{NotoSansCuneiform-Regular.ttf}[Path=\ScriptsDebugFontRoot]
\newcommand{\CU}[1]{{\CUfont #1}}

% 埃及聖書文（建議放 hyperref 前）
\usepackage{hiero}
% 建htx檔，然後終端機跑powershell -NoProfile -ExecutionPolicy Bypass -File .\build_hiero.ps1

% ---------- 蒙古文 ----------
\usepackage{graphicx}
\usepackage[export]{adjustbox}
\newfontfamily\MOVfont[
  Renderer=HarfBuzz,
  RawFeature={script=mong}
]{NotoSansMongolian-Regular.ttf}[Path=\ScriptsDebugFontRoot]
\newcommand{\MOV}[2][0pt]{%
  \raisebox{#1}{\adjustbox{valign=c}{\rotatebox[origin=c]{-90}{{\MOVfont #2}}}}%
}
